% ============================================================
% results_draft.tex
% LLM-Enhanced Dynamic Graph Networks for CVR Prediction
% Liu Yize | UCL MSc KIDS
% Draft (2026-08-08) — restructured Chapter 4 (Results) content.
%
% Organised by dataset (4.1 Ali-CCP, 4.2 Amazon), each with the same
% internal shape: main model-variant comparison first, then the
% dataset-specific extension study (Ali-CCP: encoder capacity;
% Amazon: Dynamic Graph). V3 (hybrid routing) appears under both,
% since it was run on both datasets with opposite findings.
%
% Numbers here are RESULTS ONLY — model/router/architecture
% definitions live in methods_draft.tex (\S3.2) and are referenced,
% not repeated. Tables are deliberately split (one table per
% comparison) rather than one large combined table, since the
% multi-seed mean+-std and significance columns don't fit/read well
% together at this many rows.
%
% Cite with \cite{...} using keys from references.bib (unchanged).
% All figures cross-checked against results/multiseed_comparison_summary.md,
% results/significance_tests_full.log, results/v4_significance_tests.log,
% results/v4_ablation_significance.log (2026-08-08).
% ============================================================

\chapter{Results}
\label{ch:results}

\section{Ali-CCP}
\label{sec:results-aliccp}

\subsection{Model Variant Comparison}
\label{sec:results-aliccp-main}

\begin{table}[h]
\centering
\caption{Ali-CCP: test CTCVR-AUC by model variant (MiniLM encoder), mean $\pm$ std across 5 seeds}
\label{tab:aliccp-main}
\begin{tabular}{lrrr}
\toprule
Model & Overall & Seen items & Cold-start items \\
\midrule
Baseline (ID)  & 0.5562 $\pm$ 0.0048 & 0.5572 $\pm$ 0.0105 & 0.5573 $\pm$ 0.0030 \\
V1 (item)      & 0.5576 $\pm$ 0.0248 & 0.5655 $\pm$ 0.0372 & 0.5462 $\pm$ 0.0127 \\
V1-Full (item+user) & 0.5121 $\pm$ 0.0011 & 0.4908 $\pm$ 0.0019 & 0.5421 $\pm$ 0.0031 \\
V2 (ID+align)  & 0.5574 $\pm$ 0.0050 & 0.5566 $\pm$ 0.0052 & \textbf{0.5639} $\pm$ 0.0089 \\
\bottomrule
\end{tabular}
\end{table}

\begin{table}[h]
\centering
\caption{Ali-CCP: paired significance tests, key comparisons (5 shared seeds)}
\label{tab:aliccp-sig}
\begin{tabular}{llrrl}
\toprule
Comparison & Segment & Mean diff & $p$ & Significant? \\
\midrule
Baseline vs.\ V1 & Overall & +0.0014 & 0.918 & No \\
Baseline vs.\ V2 & Overall & +0.0011 & 0.686 & No \\
V1 vs.\ V2       & Overall & --0.0002 & 0.986 & No \\
Baseline vs.\ V1 & Cold-start & --0.0111 & 0.101 & No \\
Baseline vs.\ V2 & Cold-start & +0.0065 & 0.138 & No \\
\textbf{V1 vs.\ V2} & \textbf{Cold-start} & \textbf{+0.0177} & \textbf{0.0048} & \textbf{Yes} \\
Baseline vs.\ V1 & Seen & +0.0083 & 0.702 & No \\
Baseline vs.\ V2 & Seen & --0.0007 & 0.822 & No \\
\bottomrule
\end{tabular}
\end{table}

Across 5 seeds, the only statistically significant difference among Baseline/V1/V2 on Ali-CCP is \textbf{V2 beating V1 on the cold-start segment} ($p=0.0048$). None of the Baseline-vs-V1, Baseline-vs-V2, or V1-vs-V2 overall/seen comparisons clear $\alpha=0.05$; at $n=5$ seeds these should be read as directional, not established, differences. V1-Full is omitted from the significance table (only compared descriptively here) because it is dominated by both other REPLACE-family results on every segment (Table~\ref{tab:aliccp-main}) and was not carried forward past this point.

Three qualitative patterns are nonetheless consistent with the single-seed pattern reported in earlier project stages, now confirmed not to be a seed artefact for the two REPLACE variants overall: (1) V1's cold-start AUC (0.5462) trails V2's (0.5639), consistent with the significant V1-vs-V2 cold-start gap above; (2) V1-Full is uniformly the weakest model on every segment, consistent with the low cardinality of Ali-CCP's available demographic pseudo-text fields (as few as 2--4 distinct values per field) causing many users to collapse onto near-identical text; (3) V2 is the only LLM-integrated variant that is never the weakest model on any segment, consistent with retaining learned ID embeddings (Section~\ref{sec:models}) preserving memorisation capacity that both REPLACE variants (V1, V1-Full) trade away.

\subsection{Encoder Capacity: MiniLM vs.\ MPNet}
\label{sec:results-encoders}

\begin{table}[h]
\centering
\caption{Ali-CCP test CTCVR-AUC: MiniLM (384-dim) vs.\ MPNet (768-dim), mean $\pm$ std across 5 seeds}
\label{tab:mpnet-vs-minilm}
\begin{tabular}{lrrr}
\toprule
Model & Overall & Seen items & Cold-start items \\
\midrule
V1 -- MiniLM      & 0.5576 $\pm$ 0.0248 & 0.5655 $\pm$ 0.0372 & 0.5462 $\pm$ 0.0127 \\
V1 -- MPNet       & 0.5462 $\pm$ 0.0133 & 0.5427 $\pm$ 0.0188 & 0.5499 $\pm$ 0.0126 \\
V1-Full -- MiniLM & 0.5121 $\pm$ 0.0011 & 0.4908 $\pm$ 0.0019 & 0.5421 $\pm$ 0.0031 \\
V1-Full -- MPNet  & 0.5111 $\pm$ 0.0093 & 0.5051 $\pm$ 0.0033 & 0.5187 $\pm$ 0.0243 \\
V2 -- MiniLM      & 0.5574 $\pm$ 0.0050 & 0.5566 $\pm$ 0.0052 & 0.5639 $\pm$ 0.0089 \\
V2 -- MPNet       & 0.5576 $\pm$ 0.0134 & 0.5646 $\pm$ 0.0109 & 0.5492 $\pm$ 0.0247 \\
\bottomrule
\end{tabular}
\end{table}

None of the three MiniLM-vs-MPNet overall-AUC differences is statistically significant across 5 seeds (V1: $p=0.260$; V1-Full: $p=0.820$; V2: $p=0.976$; significance was tested on the overall metric only). This is an important revision relative to this project's earlier, single-seed ($\text{seed}=42$) observation, which showed a much larger apparent MPNet cold-start regression for V2 specifically (0.5361 vs.\ 0.5724, a $-0.036$ gap) and was provisionally read as evidence that a higher-dimensional embedding space is harder to align under an unchanged $\lambda_{\text{align}}$ and epoch budget. Under 5-seed averaging that gap shrinks to $-0.0147$ (0.5492 vs.\ 0.5639, Table~\ref{tab:aliccp-main} vs.\ Table~\ref{tab:mpnet-vs-minilm}) and is not distinguishable from seed noise. The corrected conclusion is weaker but still directionally consistent: a larger, generally stronger general-purpose encoder does not reliably improve --- and numerically tends to reduce, though not significantly --- CTCVR-AUC on Ali-CCP, on both the overall and cold-start segments, which is inconsistent with ``insufficient encoder capacity'' as the explanation for the LLM-integrated variants' modest gains over Baseline (Section~\ref{sec:discussion}).

\subsection{Hybrid Segment Routing (V3)}
\label{sec:results-v3-aliccp}

Segmenting the test set by context length (Section~\ref{sec:context-segments}) crossed with seen/cold-start status reveals a crossover the aggregate CTCVR-AUC numbers hide: V2 is the strongest model overall (Table~\ref{tab:aliccp-main}), but in the single hardest cell --- short-context \& cold-start (418{,}643 rows, 20.9\% of test) --- V1 (0.5671) clearly beats V2 (0.5455) and Baseline (0.5413), exactly as the supervision-asymmetry argument in Section~\ref{sec:v3} predicts.

\begin{table}[h]
\centering
\caption{Ali-CCP test CTCVR-AUC: V3-Hybrid vs.\ its four component models (single seed, $\text{seed}=42$)}
\label{tab:v3-results-aliccp}
\begin{tabular}{lrr}
\toprule
Model & Overall & Short-context \& cold-start \\
\midrule
Baseline    & 0.5610 & 0.5413 \\
V1          & 0.5299 & 0.5671 \\
V1-Full     & 0.5124 & 0.5564 \\
V2          & 0.5616 & 0.5455 \\
V3-Hybrid   & \textbf{0.5649} & \textbf{0.5671} \\
\bottomrule
\end{tabular}
\end{table}

V3-Hybrid wins on both counts, not only the hard segment it targets. It exactly matches V1's number in that cell (0.5671, as expected --- it is literally V1's prediction there) while also improving \emph{overall} CTCVR-AUC over pure V2 (0.5649 vs.\ 0.5616), despite V1 being the substantially weaker model everywhere else and routed to only 20.9\% of test rows. This was not guaranteed a priori (Section~\ref{sec:v3}): combining two independently-calibrated models' raw probability outputs risked hurting the pooled ranking even if each branch were locally correct. That it did not is an empirical property of this model pair on this test set. This result is single-seed only (Section~\ref{sec:metrics-multiseed}) and has not been significance-tested.

\section{Amazon Reviews'23}
\label{sec:results-amazon}

\subsection{Model Variant Comparison}
\label{sec:results-amazon-main}

\begin{table}[h]
\centering
\caption{Amazon Reviews'23 (Video\_Games) test AUC by model variant (MiniLM), mean $\pm$ std across 5 seeds}
\label{tab:amazon-main}
\begin{tabular}{lrrr}
\toprule
Model & Overall\footnotemark & Seen items & Cold-start items \\
\midrule
Baseline (ID)                & 0.4687 $\pm$ 0.0048 & 0.6150 $\pm$ 0.0054 & 0.5011 $\pm$ 0.0062 \\
V1 (item, text-replace)      & 0.4178 $\pm$ 0.0025 & 0.6047 $\pm$ 0.0064 & 0.4945 $\pm$ 0.0035 \\
V2 (ID+align)                & \textbf{0.4930} $\pm$ 0.0098 & \textbf{0.6177} $\pm$ 0.0119 & \textbf{0.5027} $\pm$ 0.0106 \\
\bottomrule
\end{tabular}
\end{table}
\footnotetext{\texttt{Test overall} pools the seen-item segment (11.8\% positive rate) and the cold-start segment (46.7\% positive rate) on the same test set. AUC does not decompose linearly across strata with very different positive rates, so \texttt{overall} can and does move in ways the segment-level numbers alone would not predict (Section~\ref{sec:discussion}); the seen/cold-start columns are the primary comparison for Amazon, not \texttt{overall}.}

\begin{table}[h]
\centering
\caption{Amazon: paired significance tests (5 shared seeds)}
\label{tab:amazon-sig}
\begin{tabular}{llrrl}
\toprule
Comparison & Segment & Mean diff & $p$ & Significant? \\
\midrule
Baseline vs.\ V1 & Overall\footnotemark\label{fn:amazon-sig-overall} & --0.0509 & $<$0.0001 & Yes \\
Baseline vs.\ V2 & Overall\footnotemark[\ref{fn:amazon-sig-overall}] & +0.0243 & 0.0036 & Yes \\
V1 vs.\ V2       & Overall\footnotemark[\ref{fn:amazon-sig-overall}] & +0.0752 & $<$0.0001 & Yes \\
Baseline vs.\ V1 & Cold-start & --0.0066 & 0.172 & No \\
Baseline vs.\ V2 & Cold-start & +0.0016 & 0.589 & No \\
V1 vs.\ V2       & Cold-start & +0.0082 & 0.257 & No \\
Baseline vs.\ V2 & Seen & +0.0027 & 0.676 & No \\
\bottomrule
\end{tabular}
\end{table}
\footnotetext{Significant on the confounded pooled metric only (see previous footnote) --- since none of the corresponding seen/cold-start subgroup differences are themselves significant, these overall-AUC gaps most likely reflect each model's score distribution interacting with the segment positive-rate mismatch, not a genuine ranking-quality difference. Not cited as evidence of a real Baseline/V1/V2 gap in Section~\ref{sec:discussion}.}

Both Baseline and V2 learn real signal on seen items ($\sim$0.61--0.62 AUC, well above chance); all three models are close to chance on cold-start items ($\sim$0.49--0.50), and none of the pairwise cold-start differences reach significance. V1 (REPLACE) is the weakest model on every segment, reproducing --- now with 5-seed confirmation rather than a single run --- the Ali-CCP finding that naively substituting a frozen text embedding for a learned item ID embedding underperforms keeping the ID embedding, even with genuine natural-language item text (Section~\ref{sec:discussion}).

\subsection{Encoder Capacity: MiniLM vs.\ MPNet}
\label{sec:results-amazon-encoders}

The same MiniLM-vs-MPNet swap tested on Ali-CCP (Section~\ref{sec:results-encoders}) was repeated on Amazon for the two variants that admit it, V1 (REPLACE) and V2 (ALIGN), each at 5 seeds.

\begin{table}[h]
\centering
\caption{Amazon test AUC: MiniLM (384-dim) vs.\ MPNet (768-dim), mean $\pm$ std across 5 seeds}
\label{tab:amazon-mpnet-vs-minilm}
\begin{tabular}{lrrr}
\toprule
Model & Overall\footnotemark & Seen items & Cold-start items \\
\midrule
V1 -- MiniLM & 0.4178 $\pm$ 0.0025 & 0.6047 $\pm$ 0.0064 & 0.4945 $\pm$ 0.0035 \\
V1 -- MPNet  & 0.4721 $\pm$ 0.0074 & 0.5815 $\pm$ 0.0190 & 0.5331 $\pm$ 0.0078 \\
V2 -- MiniLM & 0.4930 $\pm$ 0.0098 & 0.6177 $\pm$ 0.0119 & 0.5027 $\pm$ 0.0106 \\
V2 -- MPNet  & 0.4825 $\pm$ 0.0047 & 0.6196 $\pm$ 0.0044 & 0.5033 $\pm$ 0.0123 \\
\bottomrule
\end{tabular}
\end{table}
\footnotetext{See Table~\ref{tab:amazon-main}'s footnote: \texttt{overall} pools segments with very different positive rates and is not the primary comparison.}

\begin{table}[h]
\centering
\caption{Amazon MiniLM vs.\ MPNet: paired significance tests (5 shared seeds)}
\label{tab:amazon-encoder-sig}
\begin{tabular}{llrrl}
\toprule
Model & Segment & Mean diff (MPNet $-$ MiniLM) & $p$ & Significant? \\
\midrule
V1 & Overall     & +0.0542 & 0.0002 & Yes \\
V1 & Seen        & --0.0232 & 0.0225 & Yes (MPNet worse) \\
V1 & Cold-start  & +0.0386 & 0.0001 & Yes \\
V2 & Overall     & --0.0105 & 0.089 & No \\
V2 & Seen        & +0.0019 & 0.749 & No \\
V2 & Cold-start  & +0.0006 & 0.953 & No \\
\bottomrule
\end{tabular}
\end{table}

Unlike Ali-CCP, where no MiniLM-vs-MPNet difference reached significance for any variant (Section~\ref{sec:results-encoders}), the encoder swap has a real, replicable effect on Amazon's V1 (REPLACE) --- but not a uniformly positive one: MPNet significantly improves both overall AUC ($p=0.0002$) and, more importantly given the pooling caveat, the cold-start segment ($p=0.0001$), while significantly \emph{hurting} the seen-items segment ($p=0.0225$). A plausible reading follows from V1's design (Section~\ref{sec:models}): V1 has no learned ID embedding at all, so its item representation is entirely the frozen text embedding passed through a small trainable adapter --- the encoder's raw representational quality is not a secondary signal here, it is the only signal. MPNet's larger, generally stronger general-purpose embedding space transfers better to items the adapter never trained on directly (cold-start), consistent with the cold-start gain; but with a fixed adapter capacity and epoch budget, that same higher-dimensional space is apparently harder to compress into a maximally discriminative representation for items the model \emph{does} get to memorise (seen), consistent with the seen-item regression. This is a capacity/compression trade-off, not a strict improvement.

V2 (ALIGN) shows no significant difference on any segment ($p=0.089$--$0.953$), matching the Ali-CCP V2 result exactly. The likely explanation is the same one given there and in Section~\ref{sec:v3}: V2's text embedding only ever receives gradient through the contrastive alignment loss, never the main CTR/CVR loss, so a better encoder has a weaker causal path to improved predictions than it does under V1's direct supervision.

Read together with Section~\ref{sec:results-encoders}, this is the only encoder-capacity comparison across both datasets that produces a significant effect, and it appears specifically for the REPLACE pattern (V1) on the dataset with genuine natural-language item text (Amazon), not for ALIGN (V2) on either dataset, nor for REPLACE on Ali-CCP's anonymised pseudo-text. This is consistent with REPLACE variants being more sensitive to encoder quality than ALIGN variants by construction --- REPLACE substitutes the text embedding as the item's \emph{only} representation, while ALIGN only recruits it for a secondary regularising loss and a cold-start fallback --- but the seen-item regression means this sensitivity is not simply "bigger encoder, better model"; at $n=5$ seeds the practical takeaway is that encoder choice is a real design decision for REPLACE-style integration on datasets with genuine text, not a free upgrade.

\subsection{Hybrid Segment Routing (V3)}
\label{sec:results-v3-amazon}

The router defined in Section~\ref{sec:v3} was first run on Amazon with Ali-CCP's routing direction unchanged (hard segment $\to$ V1). This is the wrong direction for Amazon: a per-cell breakdown showed \textbf{V2} wins the hard segment here (0.5173 vs.\ V1's 0.4872), while V1 wins the other three cells (long \& seen, long \& cold-start, short \& seen) --- the reverse of Ali-CCP, where V1 wins the hard segment (Section~\ref{sec:results-v3-aliccp}). The rule was flipped (hard segment $\to$ V2, elsewhere $\to$ V1) and rerun.

\begin{table}[h]
\centering
\caption{Amazon test AUC: V3-Hybrid before and after the routing-direction fix (single seed, $\text{seed}=42$)}
\label{tab:v3-results-amazon}
\begin{tabular}{lrr}
\toprule
Model & Overall & Hard segment (short-context \& cold-start) \\
\midrule
Baseline (ID)                       & 0.4698 & n/a \\
V1 alone                            & 0.4209 & 0.4872 \\
V2 alone                            & 0.5088 & 0.5173 \\
V3-Hybrid (Ali-CCP direction, wrong for Amazon) & 0.4443 & 0.4872 \\
V3-Hybrid (corrected direction)     & \textbf{0.4944} & \textbf{0.5173} \\
\bottomrule
\end{tabular}
\end{table}

This flip is itself a finding, not just a bug fix: it is the direct empirical answer to whether the identity of the ``right'' branch for the hard segment (Section~\ref{sec:v3}) is a fixed property of the router design or depends on the underlying text. It does not transfer --- on Ali-CCP's pseudo-text, V1 (direct task supervision on anonymised template text) wins the hard segment; on Amazon's genuine natural-language text, V2 (ID memorisation plus a weakly-supervised alignment projection) wins it instead. The corrected hybrid matches or beats V2 in every individual cell (V1's cells: long \& seen 0.6061 vs.\ V2's 0.5970, long \& cold 0.4994 vs.\ V2's 0.4877, short \& seen 0.6070 vs.\ V2's 0.5988; the hard-segment cell ties V2 exactly at 0.5173) and beats Baseline overall (0.4944 vs.\ 0.4698). It does \emph{not}, however, beat V2 alone on pooled \texttt{overall} AUC (0.4944 vs.\ 0.5088), despite winning or tying every individual cell against it --- the same pooling effect noted in Table~\ref{tab:amazon-main}'s footnote (hard-segment cells run 46--47\% positive rate vs.\ 12\% for seen cells), not a routing error. Unlike Ali-CCP, where the hybrid's overall AUC (0.5649, Table~\ref{tab:v3-results-aliccp}) cleanly beat every individual model, on Amazon the per-cell breakdown is the number to cite as the router's headline result, not \texttt{overall}. This result is single-seed only and has not been significance-tested; the original, mis-routed run is preserved in the project repository for reference.

\subsection{Dynamic Graph Layer (V4)}
\label{sec:results-v4}

\begin{table}[h]
\centering
\caption{Amazon test AUC: Dynamic Graph layer (V4) and its ID-only ablation, mean $\pm$ std across 5 seeds}
\label{tab:v4-results}
\begin{tabular}{lrrr}
\toprule
Model & Overall & Seen items & Cold-start items \\
\midrule
Baseline (ID, no graph)              & 0.4687 $\pm$ 0.0048 & 0.6150 $\pm$ 0.0054 & 0.5011 $\pm$ 0.0062 \\
V2 (ID + LLM align, no graph)        & 0.4930 $\pm$ 0.0098 & 0.6177 $\pm$ 0.0119 & 0.5027 $\pm$ 0.0106 \\
V4-ID-Only (ID + graph, no LLM)      & 0.4637 $\pm$ 0.0095 & 0.6091 $\pm$ 0.0079 & 0.5010 $\pm$ 0.0038 \\
V4 (ID + LLM align + graph)          & 0.4807 $\pm$ 0.0085 & 0.6207 $\pm$ 0.0129 & 0.5014 $\pm$ 0.0117 \\
\bottomrule
\end{tabular}
\end{table}

\begin{table}[h]
\centering
\caption{V4: paired significance tests (5 shared seeds)}
\label{tab:v4-sig}
\begin{tabular}{llrrl}
\toprule
Comparison & Segment & Mean diff & $p$ & Significant? \\
\midrule
Baseline vs.\ V4-ID-Only & Overall & --0.0051 & 0.396 & No \\
Baseline vs.\ V4-ID-Only & Seen & --0.0058 & 0.341 & No \\
Baseline vs.\ V4-ID-Only & Cold-start & --0.0000 & 0.995 & No \\
V4-ID-Only vs.\ V4 & Overall & +0.0171 & 0.061 & No (borderline) \\
V4-ID-Only vs.\ V4 & Cold-start & +0.0004 & 0.938 & No \\
\textbf{V2 vs.\ V4} & \textbf{Overall} & \textbf{--0.0123} & \textbf{0.022} & \textbf{Yes (V4 worse)} \\
V2 vs.\ V4 & Cold-start & --0.0013 & 0.885 & No \\
V2 vs.\ V4 & Seen & +0.0030 & 0.613 & No \\
Baseline vs.\ V2 (reference, from Table~\ref{tab:amazon-sig}) & Overall & +0.0243 & 0.0036 & Yes \\
\bottomrule
\end{tabular}
\end{table}

This is a clean ablation of the two layers V4 combines. The graph mechanism alone, with no LLM signal (Baseline vs.\ V4-ID-Only), produces no detectable improvement over the plain ID baseline on any segment, including cold-start ($p=0.995$) --- the segment this mechanism was specifically expected to help most (Section~\ref{sec:metrics}). Adding LLM alignment on top of the graph mechanism (V4-ID-Only vs.\ V4) trends positive on overall AUC but does not clear significance ($p=0.061$, borderline at $n=5$). Adding the graph mechanism on top of an already-working LLM-alignment model (V2 vs.\ V4) makes things significantly \emph{worse} on overall AUC ($p=0.022$) and has no effect on cold-start. For reference, V2's gain over Baseline \emph{is} significant (bottom row) --- the LLM-alignment layer has a real, replicable effect that the graph aggregator, on this dataset, does not.

The graph mechanism is not simply broken: its learned recency-decay rate converges to $\approx 0.6$ (not collapsed to zero), indicating the attention mechanism is actively using the $\Delta t$ signal rather than ignoring it. The most likely explanation is data sparsity rather than a mechanism failure: per-user history at prediction time has a median of only 1--2 available items among rows that have any history at all, 18.9\% of test rows have zero prior history (a user's first-ever interaction), and 0.0\% of users ever reach the 50-item history cap (Section~\ref{sec:v4}) --- there is little sequential structure in this dataset's density regime for a target-attention aggregator to exploit, regardless of how well the mechanism itself is trained. This is treated as a genuine, reportable negative result for the Dynamic Graph layer specifically, in contrast to the LLM Encoder layer (Section~\ref{sec:results-amazon-main}), not as a target for further tuning against, given the project's time budget (Section~\ref{sec:discussion}).
